\section{Introdução}\label{sec:introducao}

A resolução de labirintos é uma tarefa clássica que demanda raciocínio espacial, planejamento e adaptação a ambientes estruturados. Tradicionalmente, algoritmos simbólicos e métodos baseados em busca foram empregados para solucionar esse tipo de problema. No entanto, com o avanço dos Modelos de Linguagem de Grande Escala (\textit{Large Language Models} - LLMs), surge a oportunidade de investigar como essas arquiteturas, originalmente projetadas para linguagem natural, podem ser adaptadas para tarefas que exigem compreensão espacial e manipulação de representações simbólicas.

Apesar do sucesso dos LLMs em tarefas linguísticas, ainda há lacunas no entendimento de como esses modelos internalizam e representam informações espaciais, especialmente após processos de especialização como o \textit{fine-tuning}. O caráter de ``caixa-preta'' dos LLMs dificulta a análise dos mecanismos internos responsáveis pelo aprendizado e pela tomada de decisão, tornando essencial o uso de técnicas de explicabilidade para investigar as transformações ocorridas nas redes neurais durante o treinamento para tarefas específicas.

Neste contexto, este trabalho propõe-se a investigar como LLMs abertos, neste estudo com 1,5 bilhões de parâmetros, aprendem a resolver labirintos representados em ASCII por meio de diferentes estratégias de \textit{fine-tuning}. O objetivo central é analisar as mudanças nas ativações internas das redes neurais associadas ao aprendizado da tarefa, utilizando a biblioteca LLM-MRI \cite{costa2024llmmri} para visualizar e comparar as representações neuronais antes e depois do processo de aprendizagem além de avaliar métricas de redes complexas aplicadas aos grafos de ativações gerados.

Os objetivos específicos deste estudo são: (i) verificar como diferentes estratégias de \textit{fine-tuning} refletem nas ativações internas dos modelos; (ii) analisar as estruturas dos grafos de ativações recorrentes gerados a partir dos embeddings dos modelos; (iii) avaliar métricas de grafos para identificar padrões emergentes e alterações nas representações internas; e (iv) discutir as implicações dos resultados para o entendimento do raciocínio espacial em LLMs.

Ao abordar essas questões, este trabalho visa contribuir tanto para o avanço no entendimento prático das técnicas de especialização de LLMs quanto para o aprofundamento teórico sobre os processos de representação e raciocínio em redes neurais profundas. O foco está no entendimento dos mecanismos de raciocínio espacial em LLMs, avaliando desempenho e interpretabilidade dos modelos em tarefas estruturadas de navegabilidade em labirintos textuais.

O restante desta proposta está organizado da seguinte forma: a Seção \ref{sec:fund-trabs-rel} evidencia os trabalhos relacionados e os fundamentos teóricos necessários para o entendimento do tema; a Seção \ref{sec:metodologia} detalha a metodologia que foi adotada para o desenvolvimento deste estudo, incluindo a preparação dos dados, seleção de modelos, estratégias de ajuste e análise das métricas sobre os grafos de ativações neurais; os resultados obtidos são discutidos na Seção \ref{sec:resultados}; por fim, a Seção \ref{sec:discussao} apresenta as conclusões, limitações e perspectivas para o aprofundamento do estudo.